\documentclass[12pt,a4paper]{article}

% ===== PACKAGES =====
\usepackage[utf8]{inputenc}
\usepackage[T5]{fontenc}
\usepackage[vietnamese,english]{babel}
\usepackage{geometry}
\usepackage{graphicx}
\usepackage{hyperref}
\usepackage{listings}
\usepackage{xcolor}
\usepackage{booktabs}
\usepackage{array}
\usepackage{longtable}
\usepackage{tabularx}
\usepackage{enumitem}
\usepackage{titlesec}
\usepackage{fancyhdr}
\usepackage{amsmath}
\usepackage{float}
\usepackage{caption}
\usepackage{subcaption}
\usepackage{tikz}
\usetikzlibrary{shapes.geometric, arrows.meta, positioning, fit, backgrounds,
                decorations.pathmorphing, decorations.markings}
\usepackage{mdframed}
\usepackage{tcolorbox}
\tcbuselibrary{skins, breakable}
\usepackage{pgfornament}  % for corner ornaments (optional, graceful fallback)

% ===== PAGE GEOMETRY =====
\geometry{
  top=2.5cm,
  bottom=2.5cm,
  left=3cm,
  right=2cm,
}

% ===== COLORS =====
\definecolor{primaryblue}{RGB}{0, 84, 166}
\definecolor{secondaryblue}{RGB}{70, 130, 180}
\definecolor{lightgray}{RGB}{245, 245, 245}
\definecolor{codebg}{RGB}{248, 248, 248}
\definecolor{codeframe}{RGB}{180, 180, 180}
\definecolor{keywordcolor}{RGB}{0, 100, 200}
\definecolor{stringcolor}{RGB}{160, 0, 0}
\definecolor{commentcolor}{RGB}{100, 100, 100}
\definecolor{successgreen}{RGB}{0, 128, 0}

% ===== HYPERREF SETUP =====
\hypersetup{
  colorlinks=true,
  linkcolor=primaryblue,
  urlcolor=secondaryblue,
  citecolor=primaryblue,
  pdftitle={Assignment 05 - BookStore Microservice System},
  pdfauthor={SA\&D 2026},
}

% ===== LISTINGS SETUP =====
\lstset{
  backgroundcolor=\color{codebg},
  basicstyle=\ttfamily\small,
  breakatwhitespace=false,
  breaklines=true,
  captionpos=b,
  commentstyle=\color{commentcolor}\itshape,
  frame=single,
  rulecolor=\color{codeframe},
  keepspaces=true,
  keywordstyle=\color{keywordcolor}\bfseries,
  numbers=left,
  numbersep=8pt,
  numberstyle=\tiny\color{gray},
  showspaces=false,
  showstringspaces=false,
  showtabs=false,
  stepnumber=1,
  stringstyle=\color{stringcolor},
  tabsize=4,
  xleftmargin=15pt,
}

\lstdefinelanguage{YAML}{
  keywords={true,false,null,yes,no},
  keywordstyle=\color{keywordcolor}\bfseries,
  sensitive=false,
  comment=[l]{\#},
  commentstyle=\color{commentcolor}\itshape,
  stringstyle=\color{stringcolor},
  morestring=[b]',
  morestring=[b]",
}

% ===== HEADER / FOOTER =====
\pagestyle{fancy}
\fancyhf{}
\fancyhead[L]{\small\textcolor{primaryblue}{\textbf{SA\&D 2026 -- Assignment 05}}}
\fancyhead[R]{\small\textcolor{gray}{BookStore Microservice System}}
\fancyfoot[C]{\thepage}
\renewcommand{\headrulewidth}{0.4pt}
\renewcommand{\footrulewidth}{0pt}

% ===== SECTION FORMATTING =====
\titleformat{\section}
  {\Large\bfseries\color{primaryblue}}
  {\thesection}{1em}{}[\titlerule]
\titleformat{\subsection}
  {\large\bfseries\color{secondaryblue}}
  {\thesubsection}{1em}{}
\titleformat{\subsubsection}
  {\normalsize\bfseries}
  {\thesubsubsection}{1em}{}

% ===== TCOLORBOX STYLES =====
\tcbset{
  servicebox/.style={
    enhanced,
    colback=lightgray,
    colframe=primaryblue,
    fonttitle=\bfseries,
    attach boxed title to top left={yshift=-2mm, xshift=5mm},
    boxed title style={colback=primaryblue, colframe=primaryblue},
    breakable,
  }
}

% ===== DOCUMENT START =====
\begin{document}

% ===== TITLE PAGE (PTIT style) =====
\begin{titlepage}
\centering
\thispagestyle{empty}

% --- Outer decorative border ---
\begin{tikzpicture}[remember picture, overlay]
  % Thick outer frame
  \draw[line width=3pt, black!80]
    ([xshift=1.2cm, yshift=-1.2cm]current page.north west)
    rectangle
    ([xshift=-1.2cm, yshift=1.2cm]current page.south east);
  % Thin inner frame
  \draw[line width=1pt, black!60]
    ([xshift=1.45cm, yshift=-1.45cm]current page.north west)
    rectangle
    ([xshift=-1.45cm, yshift=1.45cm]current page.south east);

  % ---- Corner ornaments (TikZ-drawn rosettes) ----
  \foreach \cx/\cy in {
      {[xshift=1.2cm,  yshift=-1.2cm]current page.north west} /,
      {[xshift=-1.2cm, yshift=-1.2cm]current page.north east} /,
      {[xshift=1.2cm,  yshift=1.2cm]current page.south west} /,
      {[xshift=-1.2cm, yshift=1.2cm]current page.south east} /
  }{} % placeholder — real ornaments drawn below

  % Top-left corner flower
  \begin{scope}[shift={([xshift=1.2cm,yshift=-1.2cm]current page.north west)}]
    \foreach \a in {0,45,...,315}
      \fill[black!70, rotate=\a] (0.18,0) ellipse (0.18 and 0.07);
    \fill[black!80] (0,0) circle (0.1);
  \end{scope}
  % Top-right corner flower
  \begin{scope}[shift={([xshift=-1.2cm,yshift=-1.2cm]current page.north east)}]
    \foreach \a in {0,45,...,315}
      \fill[black!70, rotate=\a] (0.18,0) ellipse (0.18 and 0.07);
    \fill[black!80] (0,0) circle (0.1);
  \end{scope}
  % Bottom-left corner flower
  \begin{scope}[shift={([xshift=1.2cm,yshift=1.2cm]current page.south west)}]
    \foreach \a in {0,45,...,315}
      \fill[black!70, rotate=\a] (0.18,0) ellipse (0.18 and 0.07);
    \fill[black!80] (0,0) circle (0.1);
  \end{scope}
  % Bottom-right corner flower
  \begin{scope}[shift={([xshift=-1.2cm,yshift=1.2cm]current page.south east)}]
    \foreach \a in {0,45,...,315}
      \fill[black!70, rotate=\a] (0.18,0) ellipse (0.18 and 0.07);
    \fill[black!80] (0,0) circle (0.1);
  \end{scope}
\end{tikzpicture}

% --- Page content ---
\vspace*{0.5cm}

{\large\bfseries HỌC VIỆN CÔNG NGHỆ BƯU CHÍNH VIỄN THÔNG}\\[0.4cm]
{\normalsize\bfseries KHOA CÔNG NGHỆ THÔNG TIN 1}\\[0.6cm]

\rule{0.25\textwidth}{0.8pt}\\[1.2cm]

% --- PTIT Logo (TikZ) ---
\begin{tikzpicture}[scale=1.3]
  % Outer thin circle
  \draw[line width=1.2pt, black!80] (0,0) circle (1.5cm);
  % Red filled arc (left side ~270 deg)
  \fill[red!80!black] (0,0) -- (90:1.42cm) arc[start angle=90, end angle=-200, radius=1.42cm] -- cycle;
  % White arc gap (right side)
  \fill[white] (0,0) -- (-200:1.42cm) arc[start angle=-200, end angle=90, radius=1.42cm] -- cycle;
  % Inner circle white cutout
  \fill[white] (0,0) circle (0.78cm);
  % Inner red ring
  \fill[red!75!black] (0,0) circle (0.78cm);
  \fill[white] (0,0) circle (0.62cm);
  % "PT" text in red circle
  \node[font=\bfseries\large, text=red!70!black] at (-0.15, 0.0) {PT};
  % "i" with yellow dot
  \node[font=\bfseries\large, text=black] at (0.38, 0.0) {i};
  \fill[yellow!80!orange] (0.38, 0.38) circle (0.08cm);
  % "T" underlined in red
  \node[font=\bfseries\large, text=red!70!black] at (0.65, 0.0) {T};
  \draw[red!70!black, line width=1.2pt] (-0.12,-0.32) -- (0.82,-0.32);
\end{tikzpicture}

\vspace{1.0cm}

{\LARGE\bfseries ASSIGNMENT 05}\\[0.8cm]

{\large\bfseries BỘ MÔN: KIẾN TRÚC VÀ THIẾT KẾ PHẦN MỀM}\\[1.5cm]

% --- Student info table ---
\begin{tikzpicture}
  \node[draw=black, line width=0.8pt, inner sep=0pt] (tbl) {%
    \begin{tabular}{|p{5.5cm}|p{5.5cm}|}
      \hline
      \centering Sinh viên thực hiện & Đào Ngọc Đức \\ \hline
      \centering Mã sinh viên        & B22DCCN221   \\ \hline
      \centering Lớp                 & D22CNPM03    \\ \hline
      \centering Nhóm lớp            & 03           \\ \hline
    \end{tabular}%
  };
  % small cross mark top-left (like the image)
  \node[anchor=north east] at (tbl.north west) {$\boldsymbol{+}$};
  % small square bottom-right
  \draw[line width=0.8pt] ([xshift=2pt,yshift=-2pt]tbl.south east)
        rectangle ++(0.35cm,-0.35cm);
\end{tikzpicture}

\vfill
{\small \today}

\end{titlepage}

% ===== TABLE OF CONTENTS =====
\tableofcontents
\newpage

% =============================================================
\section{Giới Thiệu Đề Tài}
% =============================================================

\subsection{Bối Cảnh và Mục Tiêu}

Assignment 05 yêu cầu sinh viên thực hiện việc \textbf{phân rã kiến trúc Monolithic} của hệ thống BookStore sang kiến trúc \textbf{Microservices}, đồng thời triển khai và kiểm thử toàn bộ hệ thống bằng Docker Compose.

Mục tiêu cụ thể của assignment:
\begin{itemize}[leftmargin=1.5cm]
  \item Phân tích và xác định các \textit{bounded contexts} trong hệ thống BookStore
  \item Thiết kế và cài đặt từng microservice độc lập sử dụng Django REST Framework
  \item Định nghĩa giao tiếp giữa các service (inter-service communication) thông qua REST API
  \item Đóng gói và triển khai toàn bộ hệ thống bằng Docker \& Docker Compose
  \item Cài đặt API Gateway làm điểm vào duy nhất cho hệ thống
  \item Xây dựng Recommender AI Service đề xuất sách dựa trên đánh giá người dùng
\end{itemize}

\subsection{Công Nghệ Sử Dụng}

\begin{table}[H]
  \centering
  \caption{Công nghệ và framework sử dụng trong hệ thống}
  \begin{tabularx}{\textwidth}{lX}
    \toprule
    \textbf{Thành Phần} & \textbf{Công Nghệ / Phiên Bản} \\
    \midrule
    Ngôn ngữ lập trình   & Python 3.11 \\
    Web Framework        & Django 4.x + Django REST Framework (DRF) \\
    Cơ sở dữ liệu        & SQLite (độc lập mỗi service) \\
    Container            & Docker + Docker Compose \\
    Giao tiếp services   & REST API (HTTP đồng bộ) \\
    API Gateway          & Django (Reverse Proxy + Web UI) \\
    \bottomrule
  \end{tabularx}
\end{table}

% =============================================================
\section{Kiến Trúc Hệ Thống}
% =============================================================

\subsection{Tổng Quan Kiến Trúc Microservices}

Hệ thống BookStore được phân rã thành \textbf{12 microservices độc lập}, mỗi service chịu trách nhiệm cho một miền nghiệp vụ cụ thể và được triển khai như một container Docker riêng biệt.

\begin{figure}[H]
  \centering
  \begin{tikzpicture}[
    service/.style={
      rectangle, rounded corners=4pt,
      draw=primaryblue, fill=primaryblue!10,
      minimum width=2.8cm, minimum height=0.8cm,
      font=\footnotesize\bfseries, align=center,
    },
    gateway/.style={
      rectangle, rounded corners=6pt,
      draw=primaryblue, fill=primaryblue,
      text=white,
      minimum width=4cm, minimum height=1cm,
      font=\small\bfseries, align=center,
    },
    client/.style={
      rectangle, rounded corners=4pt,
      draw=gray, fill=gray!10,
      minimum width=3cm, minimum height=0.8cm,
      font=\small\bfseries, align=center,
    },
    arrow/.style={-Stealth, thick, primaryblue},
    darrow/.style={-Stealth, thick, secondaryblue, dashed},
  ]
    % Client
    \node[client] (client) at (0, 5.5) {Browser / Client};

    % API Gateway
    \node[gateway] (gateway) at (0, 4) {API Gateway\\(Port 8000)};

    % Row 1
    \node[service] (cust) at (-6, 2) {customer-service\\:8001};
    \node[service] (staff) at (-3, 2) {staff-service\\:8002};
    \node[service] (mgr) at (0, 2) {manager-service\\:8003};
    \node[service] (cat) at (3, 2) {catalog-service\\:8004};
    \node[service] (book) at (6, 2) {book-service\\:8005};

    % Row 2
    \node[service] (cart) at (-5.5, 0) {cart-service\\:8006};
    \node[service] (order) at (-2, 0) {order-service\\:8007};
    \node[service] (pay) at (1.5, 0) {pay-service\\:8008};
    \node[service] (ship) at (4.5, 0) {ship-service\\:8009};

    % Row 3
    \node[service] (cmt) at (-2, -2) {comment-rate\\:8010};
    \node[service] (rec) at (2.5, -2) {recommender-ai\\:8011};

    % Arrows from client
    \draw[arrow] (client) -- (gateway);

    % Arrows from gateway
    \draw[arrow] (gateway) -- (cust);
    \draw[arrow] (gateway) -- (staff);
    \draw[arrow] (gateway) -- (mgr);
    \draw[arrow] (gateway) -- (cat);
    \draw[arrow] (gateway) -- (book);
    \draw[arrow] (gateway) -- (cart);
    \draw[arrow] (gateway) -- (order);
    \draw[arrow] (gateway) -- (cmt);
    \draw[arrow] (gateway) -- (rec);

    % Inter-service arrows
    \draw[darrow] (cust) -- (cart) node[midway, left, font=\tiny] {POST /carts/};
    \draw[darrow] (cart) -- (book) node[midway, above, font=\tiny] {GET /books/};
    \draw[darrow] (order) -- (cart) node[midway, above, font=\tiny] {GET /carts/};
    \draw[darrow] (order) -- (pay) node[midway, above, font=\tiny] {POST};
    \draw[darrow] (order) -- (ship) node[midway, above, font=\tiny] {POST};
    \draw[darrow] (cmt) -- (book) node[midway, below, font=\tiny] {};
    \draw[darrow] (rec) -- (cmt) node[midway, below, font=\tiny] {GET /reviews/};
    \draw[darrow] (rec) -- (book) node[midway, right, font=\tiny] {};

  \end{tikzpicture}
  \caption{Kiến trúc tổng thể hệ thống BookStore Microservices}
\end{figure}

\subsection{Danh Sách Các Microservices}

\begin{longtable}{lclX}
  \caption{Danh sách đầy đủ các microservices trong hệ thống} \\
  \toprule
  \textbf{STT} & \textbf{Service} & \textbf{Port} & \textbf{Chức Năng} \\
  \midrule
  \endfirsthead
  \multicolumn{4}{c}{\tablename\ \thetable{}. (tiếp theo)} \\
  \toprule
  \textbf{STT} & \textbf{Service} & \textbf{Port} & \textbf{Chức Năng} \\
  \midrule
  \endhead
  \bottomrule
  \endfoot
  1  & \texttt{api-gateway}          & 8000 & Web UI, điểm vào duy nhất, định tuyến tới tất cả services \\[0.3cm]
  2  & \texttt{customer-service}     & 8001 & Quản lý khách hàng, tự động tạo giỏ hàng khi đăng ký \\[0.3cm]
  3  & \texttt{staff-service}        & 8002 & Quản lý nhân viên (CRUD) \\[0.3cm]
  4  & \texttt{manager-service}      & 8003 & Quản lý quản lý cửa hàng (CRUD) \\[0.3cm]
  5  & \texttt{catalog-service}      & 8004 & Quản lý danh mục sách (Category) \\[0.3cm]
  6  & \texttt{book-service}         & 8005 & Quản lý sách (CRUD): tiêu đề, tác giả, giá, tồn kho \\[0.3cm]
  7  & \texttt{cart-service}         & 8006 & Quản lý giỏ hàng, xác thực sách qua book-service \\[0.3cm]
  8  & \texttt{order-service}        & 8007 & Tạo đơn hàng, kích hoạt thanh toán và giao hàng \\[0.3cm]
  9  & \texttt{pay-service}          & 8008 & Xử lý thanh toán \\[0.3cm]
  10 & \texttt{ship-service}         & 8009 & Xử lý giao hàng \\[0.3cm]
  11 & \texttt{comment-rate-service} & 8010 & Đánh giá và bình luận sách \\[0.3cm]
  12 & \texttt{recommender-ai-service} & 8011 & Gợi ý sách dựa trên đánh giá người dùng (AI/ML) \\[0.3cm]
\end{longtable}

\subsection{Giao Tiếp Giữa Các Services (Inter-Service Communication)}

\begin{table}[H]
  \centering
  \caption{Sơ đồ giao tiếp giữa các microservices}
  \begin{tabularx}{\textwidth}{llX}
    \toprule
    \textbf{Service Gọi} & \textbf{Service Được Gọi} & \textbf{Endpoint \& Mục Đích} \\
    \midrule
    customer-service     & cart-service           & \texttt{POST /api/carts/} -- Tự tạo giỏ hàng \\
    cart-service         & book-service           & \texttt{GET /api/books/\{id\}/} -- Xác thực sách tồn tại \\
    order-service        & cart-service           & \texttt{GET /api/carts/\{id\}/} -- Lấy items trong giỏ \\
    order-service        & book-service           & \texttt{GET /api/books/\{id\}/} -- Lấy giá sách \\
    order-service        & pay-service            & \texttt{POST /api/payments/} -- Kích hoạt thanh toán \\
    order-service        & ship-service           & \texttt{POST /api/shipments/} -- Kích hoạt giao hàng \\
    comment-rate-service & book-service           & \texttt{GET /api/books/\{id\}/} -- Xác thực sách \\
    recommender-ai-service & comment-rate-service & \texttt{GET /api/reviews/top-rated/} -- Lấy sách top rating \\
    recommender-ai-service & book-service         & \texttt{GET /api/books/\{id\}/} -- Lấy chi tiết sách \\
    \bottomrule
  \end{tabularx}
\end{table}

% =============================================================
\section{Cài Đặt và Triển Khai}
% =============================================================

\subsection{Cấu Trúc Thư Mục Dự Án}

\begin{lstlisting}[language=bash, caption={Cấu trúc thư mục tổng thể dự án}]
assign05/
|-- api-gateway/              # API Gateway + Web UI
|   |-- api_gateway/          # Django project config
|   |-- gateway/              # Views, URLs (proxy logic)
|   |-- templates/            # HTML templates (12 pages)
|   |-- Dockerfile
|   |-- requirements.txt
|   `-- entrypoint.sh
|-- book-service/             # Quan ly sach
|-- cart-service/             # Quan ly gio hang
|-- order-service/            # Quan ly don hang
|-- pay-service/              # Xu ly thanh toan
|-- ship-service/             # Xu ly giao hang
|-- customer-service/         # Quan ly khach hang
|-- staff-service/            # Quan ly nhan vien
|-- manager-service/          # Quan ly quan ly
|-- catalog-service/          # Danh muc sach
|-- comment-rate-service/     # Danh gia sach
|-- recommender-ai-service/   # Goi y sach
`-- docker-compose.yml        # Dinh nghia toan bo he thong
\end{lstlisting}

\subsection{Cấu Trúc Mỗi Microservice}

Mỗi microservice đều có cấu trúc chuẩn như sau:

\begin{lstlisting}[language=bash, caption={Cấu trúc chuẩn mỗi microservice (ví dụ: book-service)}]
book-service/
|-- Dockerfile          # Container image definition
|-- requirements.txt    # djangorestframework, requests
|-- entrypoint.sh       # Migrate + start server
|-- manage.py           # Django manager
|-- book_service/       # Django project settings
|   |-- settings.py
|   |-- urls.py
|   `-- wsgi.py
`-- app/                # Django application
    |-- models.py       # Data models (Book)
    |-- serializers.py  # DRF serializers
    |-- views.py        # API views (CRUD)
    |-- urls.py         # URL routing
    `-- admin.py
\end{lstlisting}

\subsection{Docker và Docker Compose}

\subsubsection{Dockerfile Chuẩn}

\begin{lstlisting}[language=Dockerfile, caption={Dockerfile chuẩn cho các microservice}]
FROM python:3.11-slim

WORKDIR /app

COPY requirements.txt .
RUN pip install --no-cache-dir -r requirements.txt

COPY . .

COPY entrypoint.sh /entrypoint.sh
RUN chmod +x /entrypoint.sh

EXPOSE 8000

ENTRYPOINT ["/entrypoint.sh"]
\end{lstlisting}

\subsubsection{Entrypoint Script}

\begin{lstlisting}[language=bash, caption={entrypoint.sh -- chạy migrate và khởi động server}]
#!/bin/bash
set -e

python manage.py makemigrations --noinput || true
python manage.py migrate --noinput

exec python manage.py runserver 0.0.0.0:8000
\end{lstlisting}

\subsubsection{Docker Compose Configuration}

\begin{lstlisting}[language=YAML, caption={docker-compose.yml -- cấu hình toàn bộ hệ thống}]
services:
  customer-service:
    build: ./customer-service
    ports:
      - "8001:8000"
    networks:
      - bookstore-network

  book-service:
    build: ./book-service
    ports:
      - "8005:8000"
    networks:
      - bookstore-network

  cart-service:
    build: ./cart-service
    ports:
      - "8006:8000"
    depends_on:
      - book-service
    networks:
      - bookstore-network

  order-service:
    build: ./order-service
    ports:
      - "8007:8000"
    depends_on:
      - cart-service
      - book-service
      - pay-service
      - ship-service
    networks:
      - bookstore-network

  recommender-ai-service:
    build: ./recommender-ai-service
    ports:
      - "8011:8000"
    depends_on:
      - book-service
      - comment-rate-service
    networks:
      - bookstore-network

  api-gateway:
    build: ./api-gateway
    ports:
      - "8000:8000"
    depends_on:
      - customer-service
      - staff-service
      - manager-service
      - catalog-service
      - book-service
      - cart-service
      - order-service
      - pay-service
      - ship-service
      - comment-rate-service
      - recommender-ai-service
    networks:
      - bookstore-network

networks:
  bookstore-network:
    driver: bridge
\end{lstlisting}

% =============================================================
\section{Thiết Kế Chi Tiết Các Microservices}
% =============================================================

\subsection{Book Service (Port 8005)}

Book Service quản lý toàn bộ thông tin sách trong hệ thống.

\subsubsection{Data Model}

\begin{lstlisting}[language=Python, caption={book-service/app/models.py}]
from django.db import models

class Book(models.Model):
    title      = models.CharField(max_length=255)
    author     = models.CharField(max_length=255)
    price      = models.DecimalField(max_digits=10, decimal_places=2)
    stock      = models.IntegerField(default=0)
    isbn       = models.CharField(max_length=13, blank=True)
    category_id = models.IntegerField(null=True, blank=True)

    def __str__(self):
        return self.title
\end{lstlisting}

\subsubsection{API Endpoints}

\begin{table}[H]
  \centering
  \caption{API endpoints của Book Service}
  \begin{tabularx}{\textwidth}{llX}
    \toprule
    \textbf{Method} & \textbf{Endpoint} & \textbf{Chức Năng} \\
    \midrule
    GET    & \texttt{/api/books/}       & Lấy danh sách tất cả sách \\
    POST   & \texttt{/api/books/}       & Thêm sách mới \\
    GET    & \texttt{/api/books/\{id\}/} & Lấy chi tiết một cuốn sách \\
    PUT    & \texttt{/api/books/\{id\}/} & Cập nhật thông tin sách \\
    DELETE & \texttt{/api/books/\{id\}/} & Xóa sách \\
    \bottomrule
  \end{tabularx}
\end{table}

\subsection{Customer Service (Port 8001)}

Customer Service xử lý việc đăng ký khách hàng và tự động tạo giỏ hàng thông qua giao tiếp với Cart Service.

\subsubsection{Luồng Xử Lý Đăng Ký Khách Hàng}

\begin{lstlisting}[language=Python, caption={Tự động tạo giỏ hàng khi đăng ký -- customer-service}]
class CustomerListCreate(APIView):
    def post(self, request):
        serializer = CustomerSerializer(data=request.data)
        if serializer.is_valid():
            customer = serializer.save()
            # Goi sang cart-service de tu dong tao gio hang
            try:
                requests.post(
                    f"{CART_SERVICE_URL}/api/carts/",
                    json={"customer_id": customer.id}
                )
            except requests.exceptions.ConnectionError:
                pass  # Cart service co the chua san sang
            return Response(serializer.data,
                          status=status.HTTP_201_CREATED)
        return Response(serializer.errors,
                      status=status.HTTP_400_BAD_REQUEST)
\end{lstlisting}

\subsection{Order Service (Port 8007)}

Order Service là service phức tạp nhất, điều phối việc tạo đơn hàng bằng cách giao tiếp với nhiều services khác.

\subsubsection{Luồng Xử Lý Đặt Hàng}

Khi khách hàng đặt hàng, luồng xử lý như sau:

\begin{enumerate}[leftmargin=1.5cm]
  \item Lấy danh sách items từ \texttt{cart-service} theo \texttt{customer\_id}
  \item Tính tổng tiền bằng cách gọi \texttt{book-service} lấy giá từng sách
  \item Tạo bản ghi \texttt{Order} và các \texttt{OrderItem} trong database
  \item Kích hoạt thanh toán bằng cách gọi \texttt{POST /api/payments/} tới \texttt{pay-service}
  \item Kích hoạt giao hàng bằng cách gọi \texttt{POST /api/shipments/} tới \texttt{ship-service}
  \item Trả về thông tin đơn hàng đã tạo
\end{enumerate}

\begin{lstlisting}[language=Python, caption={Luồng tạo đơn hàng -- order-service/app/views.py}]
CART_SERVICE_URL = "http://cart-service:8000"
PAY_SERVICE_URL  = "http://pay-service:8000"
SHIP_SERVICE_URL = "http://ship-service:8000"
BOOK_SERVICE_URL = "http://book-service:8000"

class OrderCreate(APIView):
    def post(self, request):
        customer_id      = request.data.get("customer_id")
        shipping_address = request.data.get("shipping_address", "")
        payment_method   = request.data.get("payment_method", "cash")

        # 1. Lay gio hang tu cart-service
        cart_items = []
        r = requests.get(
            f"{CART_SERVICE_URL}/api/carts/{customer_id}/")
        if r.status_code == 200:
            cart_items = r.json().get("items", [])

        # 2. Tinh tong tien qua book-service
        total = 0
        for item in cart_items:
            br = requests.get(
                f"{BOOK_SERVICE_URL}/api/books/{item['book_id']}/")
            if br.status_code == 200:
                book_price = float(br.json().get("price", 0))
                total += book_price * item["quantity"]

        # 3. Tao don hang
        order = Order.objects.create(
            customer_id=customer_id,
            total_amount=total,
            shipping_address=shipping_address,
            status='pending'
        )

        # 4. Kich hoat thanh toan
        requests.post(f"{PAY_SERVICE_URL}/api/payments/",
                     json={"order_id": order.id, "amount": str(total),
                           "method": payment_method})

        # 5. Kich hoat giao hang
        requests.post(f"{SHIP_SERVICE_URL}/api/shipments/",
                     json={"order_id": order.id,
                           "address": shipping_address})

        return Response(OrderSerializer(order).data,
                       status=status.HTTP_201_CREATED)
\end{lstlisting}

\subsection{Recommender AI Service (Port 8011)}

Recommender AI Service thực hiện gợi ý sách thông minh dựa trên đánh giá của người dùng.

\subsubsection{Chiến Lược Gợi Ý}

\begin{enumerate}[leftmargin=1.5cm]
  \item Lấy danh sách sách được đánh giá cao nhất từ \texttt{comment-rate-service}
  \item Lấy chi tiết từng cuốn sách qua \texttt{book-service}
  \item Kết hợp thông tin sách + điểm đánh giá trung bình
  \item \textbf{Fallback}: Nếu chưa có đánh giá nào, trả về 10 sách đầu tiên
\end{enumerate}

\subsubsection{Gợi Ý Cá Nhân Hóa}

\begin{lstlisting}[language=Python, caption={Gợi ý sách cá nhân hóa theo customer\_id}]
class RecommendForCustomer(APIView):
    def get(self, request, customer_id):
        # Lay tat ca danh gia
        all_reviews = requests.get(
            f"{COMMENT_RATE_SERVICE_URL}/api/reviews/").json()

        # Loc ra sach ma khach hang da danh gia (loai bo chung)
        customer_rated = {
            rev["book_id"] for rev in all_reviews
            if rev["customer_id"] == customer_id
        }

        # Lay sach top rating ma chua duoc danh gia boi customer
        recommendations = []
        for item in top_rated:
            if item["book_id"] not in customer_rated:
                book = requests.get(
                    f"{BOOK_SERVICE_URL}/api/books/{item['book_id']}/"
                ).json()
                book["avg_rating"] = item["avg_rating"]
                recommendations.append(book)

        return Response({
            "customer_id": customer_id,
            "recommendations": recommendations[:5]
        })
\end{lstlisting}

\subsection{API Gateway (Port 8000)}

API Gateway đóng vai trò là \textbf{điểm vào duy nhất} (single entry point) cho toàn bộ hệ thống, cung cấp:

\begin{itemize}[leftmargin=1.5cm]
  \item \textbf{Web UI}: Giao diện HTML để tương tác với mọi chức năng
  \item \textbf{Reverse Proxy}: Chuyển tiếp request tới services tương ứng
  \item \textbf{Service Discovery}: Kết nối với tất cả 11 services khác
\end{itemize}

\subsubsection{Danh Sách Endpoints của API Gateway}

\begin{table}[H]
  \centering
  \caption{URL routing trong API Gateway}
  \begin{tabularx}{\textwidth}{lX}
    \toprule
    \textbf{URL Pattern} & \textbf{Chức Năng} \\
    \midrule
    \texttt{/}                              & Trang chủ (index) \\
    \texttt{/books/}                        & Danh sách sách \\
    \texttt{/books/add/}                    & Thêm sách (Staff) \\
    \texttt{/customers/}                    & Danh sách khách hàng \\
    \texttt{/customers/register/}           & Đăng ký khách hàng \\
    \texttt{/customers/\{id\}/cart/}         & Giỏ hàng của khách \\
    \texttt{/customers/\{id\}/cart/add/}    & Thêm sách vào giỏ \\
    \texttt{/customers/\{id\}/order/}       & Tạo đơn hàng \\
    \texttt{/orders/}                       & Lịch sử đơn hàng \\
    \texttt{/books/\{id\}/reviews/}         & Đánh giá sách \\
    \texttt{/books/\{id\}/reviews/add/}     & Thêm đánh giá \\
    \texttt{/recommendations/}             & Gợi ý sách (AI) \\
    \bottomrule
  \end{tabularx}
\end{table}

% =============================================================
\section{Yêu Cầu Chức Năng và Cài Đặt}
% =============================================================

\subsection{Use Cases Chính}

\subsubsection{UC1: Đăng Ký Khách Hàng}

\begin{tcolorbox}[servicebox, title={Use Case: Đăng Ký Khách Hàng}]
  \textbf{Actor:} Khách hàng mới \\
  \textbf{Mô tả:} Hệ thống tạo tài khoản khách hàng và tự động khởi tạo giỏ hàng \\[0.3cm]
  \textbf{Luồng chính:}
  \begin{enumerate}
    \item Client gửi \texttt{POST /customers/register/} tới API Gateway
    \item API Gateway gọi \texttt{POST /api/customers/} tới customer-service
    \item customer-service lưu thông tin và gọi \texttt{POST /api/carts/} tới cart-service
    \item cart-service tạo giỏ hàng mới với \texttt{customer\_id}
    \item Trả về thông tin khách hàng mới tạo
  \end{enumerate}
\end{tcolorbox}

\subsubsection{UC2: Nhân Viên Quản Lý Sách}

\begin{tcolorbox}[servicebox, title={Use Case: Quản Lý Sách (Staff)}]
  \textbf{Actor:} Nhân viên (Staff) \\
  \textbf{Mô tả:} Nhân viên thực hiện CRUD sách thông qua book-service \\[0.3cm]
  \textbf{Các thao tác:}
  \begin{itemize}
    \item \textbf{Thêm sách:} \texttt{POST /api/books/} với title, author, price, stock, isbn
    \item \textbf{Xem sách:} \texttt{GET /api/books/} hoặc \texttt{GET /api/books/\{id\}/}
    \item \textbf{Cập nhật:} \texttt{PUT /api/books/\{id\}/} với thông tin mới
    \item \textbf{Xóa sách:} \texttt{DELETE /api/books/\{id\}/}
  \end{itemize}
\end{tcolorbox}

\subsubsection{UC3: Đặt Hàng}

\begin{tcolorbox}[servicebox, title={Use Case: Đặt Hàng}]
  \textbf{Actor:} Khách hàng \\
  \textbf{Tiền điều kiện:} Khách hàng đã đăng nhập và có items trong giỏ hàng \\[0.3cm]
  \textbf{Luồng chính:}
  \begin{enumerate}
    \item Khách hàng submit form tạo đơn hàng (địa chỉ, phương thức thanh toán)
    \item order-service lấy items từ cart-service và tính tổng tiền
    \item Tạo bản ghi Order + OrderItems trong database
    \item Gửi yêu cầu thanh toán tới pay-service (đồng bộ)
    \item Gửi yêu cầu giao hàng tới ship-service (đồng bộ)
    \item Trả về thông tin đơn hàng
  \end{enumerate}
\end{tcolorbox}

\subsubsection{UC4: Gợi Ý Sách (AI Recommender)}

\begin{tcolorbox}[servicebox, title={Use Case: Gợi Ý Sách}]
  \textbf{Actor:} Khách hàng \\
  \textbf{Mô tả:} Hệ thống gợi ý sách phù hợp dựa trên đánh giá cộng đồng \\[0.3cm]
  \textbf{Chiến lược:}
  \begin{itemize}
    \item \textbf{Global:} Lấy top sách được đánh giá cao nhất từ tất cả người dùng
    \item \textbf{Cá nhân hóa:} Lọc bỏ sách khách hàng đã đánh giá, gợi ý sách chưa đọc
    \item \textbf{Fallback:} Nếu chưa có đánh giá, trả về 10 sách đầu tiên trong hệ thống
  \end{itemize}
\end{tcolorbox}

% =============================================================
\section{Kết Quả và Kiểm Thử}
% =============================================================

\subsection{Khởi Động Hệ Thống}

Toàn bộ hệ thống được khởi động bằng một lệnh duy nhất:

\begin{lstlisting}[language=bash, caption={Lện khởi động hệ thống}]
# Khoi dong tat ca 12 services
docker compose up --build

# Kiem tra trang thai cac container
docker compose ps

# Xem log cua mot service cu the
docker compose logs -f book-service
\end{lstlisting}

Sau khi các container khởi động thành công, có thể truy cập:
\begin{itemize}[leftmargin=1.5cm]
  \item \textbf{Web UI:} \url{http://localhost:8000}
  \item \textbf{Book Service API:} \url{http://localhost:8005/api/books/}
  \item \textbf{Customer Service API:} \url{http://localhost:8001/api/customers/}
  \item \textbf{Recommender API:} \url{http://localhost:8011/api/recommendations/}
\end{itemize}

\subsection{Test Cases}

\subsubsection{Kiểm Thử API Bằng curl/Postman}

\begin{lstlisting}[language=bash, caption={Test API -- Tạo và lấy dữ liệu}]
# 1. Dang ky khach hang moi
curl -X POST http://localhost:8001/api/customers/ \
  -H "Content-Type: application/json" \
  -d '{"name": "Nguyen Van A", "email": "nva@example.com"}'

# 2. Them sach moi
curl -X POST http://localhost:8005/api/books/ \
  -H "Content-Type: application/json" \
  -d '{"title": "Clean Code", "author": "Robert Martin",
       "price": 150000, "stock": 10}'

# 3. Lay danh sach sach
curl http://localhost:8005/api/books/

# 4. Them sach vao gio hang (customer_id=1, cart_id=1)
curl -X POST http://localhost:8006/api/cart-items/ \
  -H "Content-Type: application/json" \
  -d '{"cart": 1, "book_id": 1, "quantity": 2}'

# 5. Tao don hang
curl -X POST http://localhost:8007/api/orders/ \
  -H "Content-Type: application/json" \
  -d '{"customer_id": 1, "shipping_address": "123 Le Loi",
       "payment_method": "cash"}'

# 6. Danh gia sach
curl -X POST http://localhost:8010/api/reviews/ \
  -H "Content-Type: application/json" \
  -d '{"book_id": 1, "customer_id": 1,
       "rating": 5, "comment": "Sach rat hay!"}'

# 7. Lay goi y sach
curl http://localhost:8011/api/recommendations/
\end{lstlisting}

\subsection{Xác Nhận Giao Tiếp Liên Dịch Vụ}

Bảng sau xác nhận các luồng giao tiếp chính giữa services đã hoạt động đúng:

\begin{table}[H]
  \centering
  \caption{Xác nhận giao tiếp giữa các services}
  \begin{tabularx}{\textwidth}{lXl}
    \toprule
    \textbf{Kịch Bản} & \textbf{Mô Tả} & \textbf{Kết Quả} \\
    \midrule
    Đăng ký KH & customer-service $\to$ cart-service tự tạo giỏ hàng & \textcolor{successgreen}{\checkmark} \\
    Thêm vào giỏ & cart-service xác thực sách qua book-service & \textcolor{successgreen}{\checkmark} \\
    Đặt hàng & order-service lấy giỏ, tính tiền, gọi pay \& ship & \textcolor{successgreen}{\checkmark} \\
    Gợi ý AI & recommender $\to$ comment-rate $\to$ book & \textcolor{successgreen}{\checkmark} \\
    API Gateway & Định tuyến đúng tới từng service & \textcolor{successgreen}{\checkmark} \\
    \bottomrule
  \end{tabularx}
\end{table}

% =============================================================
\section{Ưu Điểm và Hạn Chế}
% =============================================================

\subsection{Ưu Điểm Của Kiến Trúc Microservices}

\begin{itemize}[leftmargin=1.5cm]
  \item \textbf{Độc lập triển khai:} Mỗi service có thể được cập nhật, scale và deploy riêng biệt
  \item \textbf{Cô lập lỗi:} Khi một service gặp sự cố, các service khác vẫn hoạt động bình thường (circuit breaker tự nhiên qua \texttt{try/except ConnectionError})
  \item \textbf{Độc lập công nghệ:} Mỗi service có thể dùng ngôn ngữ/framework riêng (trong project này dùng Django đồng nhất)
  \item \textbf{Scale riêng lẻ:} Book Service có thể chạy nhiều instance hơn Order Service nếu cần
  \item \textbf{Cơ sở dữ liệu độc lập:} Mỗi service có SQLite riêng, tránh coupling ở tầng data
\end{itemize}

\subsection{Hạn Chế Hiện Tại}

\begin{itemize}[leftmargin=1.5cm]
  \item \textbf{Giao tiếp đồng bộ:} Sử dụng REST đồng bộ, nếu một service không phản hồi, order sẽ tạo thiếu thông tin
  \item \textbf{Không có authentication:} Chưa có JWT hoặc OAuth để bảo mật API
  \item \textbf{SQLite giới hạn:} SQLite không phù hợp môi trường production (nên dùng PostgreSQL)
  \item \textbf{Không có load balancing:} Chỉ có một instance của mỗi service
  \item \textbf{Thiếu health checks:} \texttt{depends\_on} trong Docker Compose không đảm bảo service đã sẵn sàng nhận request
\end{itemize}

\subsection{Hướng Phát Triển (Assignment 06)}

Assignment 06 sẽ mở rộng hệ thống theo hướng:
\begin{itemize}[leftmargin=1.5cm]
  \item Thêm \textbf{Message Queue} (RabbitMQ/Kafka) cho giao tiếp bất đồng bộ
  \item Tích hợp \textbf{Service Discovery} (Consul/Eureka) thay vì hardcode URL
  \item Bổ sung \textbf{Authentication Service} với JWT
  \item Chuyển sang \textbf{PostgreSQL} cho môi trường production
  \item Thêm \textbf{Centralized Logging} và \textbf{Monitoring} (ELK Stack, Prometheus)
\end{itemize}

% =============================================================
\section{Kết Luận}
% =============================================================

Assignment 05 đã được thực hiện thành công với việc phân rã hệ thống BookStore Monolithic thành \textbf{12 microservices độc lập}, mỗi service chịu trách nhiệm cho một miền nghiệp vụ cụ thể và hoạt động trong container Docker riêng biệt.

\subsection{Tóm Tắt Thành Tích}

\begin{itemize}[leftmargin=1.5cm]
  \item \checkmark\ Thiết kế và cài đặt 12 microservices đầy đủ chức năng
  \item \checkmark\ Cài đặt API Gateway với Web UI trực quan
  \item \checkmark\ Cài đặt giao tiếp liên dịch vụ qua REST API
  \item \checkmark\ Triển khai toàn bộ hệ thống bằng Docker Compose
  \item \checkmark\ Cài đặt Recommender AI Service với chiến lược gợi ý thông minh
  \item \checkmark\ Xử lý graceful failures trong giao tiếp giữa services
\end{itemize}

Dự án thể hiện sự hiểu biết về kiến trúc microservices, nguyên tắc \textit{Single Responsibility}, \textit{Loose Coupling} và \textit{High Cohesion} trong thiết kế phần mềm hiện đại.

% =============================================================
% REFERENCES
% =============================================================
\newpage
\section*{Tài Liệu Tham Khảo}
\addcontentsline{toc}{section}{Tài Liệu Tham Khảo}

\begin{enumerate}
  \item Django REST Framework Documentation. \url{https://www.django-rest-framework.org/}
  \item Docker Compose Documentation. \url{https://docs.docker.com/compose/}
  \item Martin Fowler. \textit{Microservices}. \url{https://martinfowler.com/articles/microservices.html}
  \item Sam Newman. \textit{Building Microservices}. O'Reilly Media, 2021.
  \item Tutorial Assignment 05 \& 06 -- SA\&D 2026 Course Material.
\end{enumerate}

\end{document}
